\documentclass[a4paper]{article}     
\usepackage{amsfonts}
\usepackage{amsmath}
\usepackage{amssymb}
\usepackage{graphicx}

\begin{document}

\noindent \large Auteur: Mateusz Ragan \\
\noindent \large Studierichting: Informatica\\
\noindent \large Oefeningenreeks 3F nr. 1.7\\

\medskip

\normalsize

$f(x) = \dfrac{x+3}{(x-1)^2}$\\

VA\\

$\lim\limits_{x \rightarrow 1+} \dfrac{4}{(0^+)^2} = +\infty$\\

$\lim\limits_{x \rightarrow1-} \dfrac{4}{(0^-)^2} = +\infty$\\

HA\\

$\lim\limits_{x \rightarrow\infty} \dfrac{x+3}{(x-1)^2} = \lim\limits_{x \rightarrow\infty} \dfrac{1}{x} = 0$\\

We hebben een HA dus er is geen SA.\\

$f(x)-A= \dfrac{x+3}{(x-1)^2}$\\

$\begin{array}{r|c c c c c c}
x 		 & -\infty & -3 & & 1 & +\infty \\ \hline
f(x)-A & - & 0 & + & |^2 & + \\
\end{array}$\\

$x\rightarrow -\infty$ dan is $f$ onder $A$, $x\rightarrow +\infty$ dan is $f$ boven $A$ 





\begin{thebibliography}{999}
%\bibitem{a} Referentie
\end{thebibliography}
\end{document} 
